\documentclass[11pt,a4paper]{article}
\usepackage[T1]{fontenc}
\usepackage[utf8]{inputenc}
\usepackage[main=italian,provide=*]{babel}
\usepackage{amsmath,amssymb}
\usepackage{geometry}
\geometry{margin=2.3cm,top=2.5cm,bottom=2.7cm}
\usepackage{booktabs}
\usepackage[table]{xcolor}
\usepackage{tcolorbox}
\tcbuselibrary{skins,breakable}
\usepackage{titlesec}
\usepackage{enumitem}
\usepackage{needspace}
\usepackage{fancyhdr}
\usepackage{hyperref}
\hypersetup{colorlinks=true,linkcolor=Sepia,citecolor=Sepia,urlcolor=Sepia}

% --- palette ---
\definecolor{inkblue}{HTML}{1B3A5C}
\definecolor{lightblue}{HTML}{EAF1F8}
\definecolor{accent}{HTML}{B0562A}
\definecolor{softgray}{HTML}{6B6B6B}
\definecolor{okgreen}{HTML}{2E6B3E}
\definecolor{okgreenbg}{HTML}{EAF5EC}
\definecolor{advisorpurple}{HTML}{5B3A7A}
\definecolor{advisorbg}{HTML}{F1ECF6}
\definecolor{notebar}{HTML}{9C6B12}

% --- titoli di sezione ---
\titleformat{\section}
  {\normalfont\Large\bfseries\color{inkblue}}
  {\thesection}{0.6em}{}[{\color{inkblue!40}\titlerule}]
\titlespacing{\section}{0pt}{0.8em}{0.4em}
\titleformat{\subsection}
  {\normalfont\large\bfseries\color{accent}}
  {}{0em}{}
\titlespacing{\subsection}{0pt}{0.5em}{0.2em}

\pagestyle{fancy}
\fancyhf{}
\renewcommand{\headrulewidth}{0.4pt}
\fancyhead[L]{\small\color{softgray}La decomposizione di Pauli per Hamiltoniane quantistiche}
\fancyhead[R]{\small\color{softgray}\thepage}
\fancyfoot[C]{\small\color{softgray}Tesi triennale, Samuele Schiavi --- Relatore: Prof.\ Alessandro Chiesa}

% --- box riutilizzabili ---
\newtcolorbox{keyeq}[1][]{
  colback=lightblue, colframe=inkblue!70, boxrule=0.6pt,
  arc=2pt, left=8pt, right=8pt, top=4pt, bottom=4pt, #1
}
\newtcolorbox{panelbox}[1]{
  colback=white, colframe=softgray!50, boxrule=0.5pt,
  arc=2pt, left=7pt, right=7pt, top=3pt, bottom=3pt,
  fonttitle=\bfseries\color{inkblue}, title={#1}
}
\newtcolorbox{okbox}[1]{
  colback=okgreenbg, colframe=okgreen!60, boxrule=0.6pt,
  arc=2pt, left=7pt, right=7pt, top=3pt, bottom=3pt,
  fonttitle=\bfseries\color{okgreen}, title={#1}
}
\newtcolorbox{advisorbox}[1]{
  colback=advisorbg, colframe=advisorpurple!60, boxrule=0.6pt,
  arc=2pt, left=7pt, right=7pt, top=4pt, bottom=4pt,
  fonttitle=\bfseries\color{advisorpurple}, title={#1}
}
% nota editoriale in stile "blockquote", per le correzioni segnalate in corsivo nell'originale
\newtcolorbox{editnote}{
  colback=white, colframe=white, boxrule=0pt,
  borderline west={2.2pt}{0pt}{notebar},
  arc=0pt, left=10pt, right=4pt, top=2pt, bottom=2pt,
  fontupper=\itshape\small\color{softgray!90!black}
}

\title{\vspace{-1.5em}\color{inkblue}La decomposizione di Pauli\\[0.1em]per Hamiltoniane quantistiche}
\author{Note di lavoro --- Tesi triennale, Samuele Schiavi\\ Relatore: Prof.\ Alessandro Chiesa}
\date{}

\begin{document}
\sloppy
\maketitle
\thispagestyle{fancy}
\vspace{-2em}

\begin{advisorbox}{Nota di verifica (A.\ Chiesa / UniPR)}
Documento rivisto per coerenza con la convenzione di 6-Applications e
Crippa et al., \emph{Magnetochemistry} \textbf{7}, 117 (2021). Le
correzioni rispetto alla versione originale sono segnalate, qui, nei
riquadri laterali in corsivo.
\end{advisorbox}

\section{La base di Pauli}

In meccanica quantistica, l'Hamiltoniana $H$ di un sistema a $N$ qubit è
una matrice hermitiana di dimensioni $2^N\times2^N$. Le matrici di Pauli
a singolo qubit, insieme alla matrice identità, formano un insieme di 4
operatori:
$$\{\sigma_0=I,\quad\sigma_1=X,\quad\sigma_2=Y,\quad\sigma_3=Z\}.$$
Per un sistema di $N$ qubit, le \textbf{stringhe di Pauli} sono tutti i
prodotti tensoriali di $N$ operatori:
$$P_{\boldsymbol\mu}=\sigma_{\mu_1}\otimes\sigma_{\mu_2}\otimes\cdots\otimes\sigma_{\mu_N},
\qquad \mu_i\in\{0,1,2,3\}.$$
Il numero totale di stringhe è $4^N$. Esse formano una \textbf{base
ortonormale completa} dello spazio vettoriale delle matrici
$2^N\times2^N$ rispetto al prodotto scalare di Hilbert--Schmidt:
$$\frac{1}{2^N}\operatorname{Tr}(P_{\boldsymbol\mu}^\dagger\,P_{\boldsymbol\nu})=\delta_{\boldsymbol\mu\boldsymbol\nu}.$$
Di conseguenza, qualsiasi hamiltoniana hermitiana $H$ ammette la
decomposizione \textbf{esatta e unica}:
\begin{keyeq}
\centering
$H = \displaystyle\sum_{\boldsymbol\mu} h_{\boldsymbol\mu}\, P_{\boldsymbol\mu}$
\end{keyeq}
con coefficienti \textbf{reali} (la hermitianità di $H$ garantisce
$h_{\boldsymbol\mu}\in\mathbb R$).

\section{Come si trovano i coefficienti}

Per isolare il peso di una specifica stringa di Pauli $P_k$ si sfrutta
l'ortogonalità:
\begin{keyeq}
\centering
$h_k = \dfrac{1}{2^N}\operatorname{Tr}(P_k\,H)$
\end{keyeq}
Il calcolo è esplicito: poiché $P_k$ è hermitiana e unitaria, $P_k^2=I$,
quindi $\operatorname{Tr}(P_k\cdot P_j)=2^N\delta_{kj}$.

\section{Convenzione di notazione Qiskit (little-endian)}

In Qiskit, le stringhe di Pauli sono scritte in ordine
\textbf{little-endian}: il carattere più a \textbf{destra} della
stringa corrisponde al qubit con indice \textbf{0}, quello più a
sinistra al qubit con indice più alto.

\begin{center}
\renewcommand{\arraystretch}{1.25}
\begin{tabular}{@{}cc@{}}
\toprule
\rowcolor{inkblue!10}
\textbf{Stringa Qiskit} & \textbf{Operatore matematico}\\
\midrule
\texttt{"ZI"} & $Z_1\otimes I_0$\\
\texttt{"IZ"} & $I_1\otimes Z_0$\\
\texttt{"XX"} & $X_1\otimes X_0$\\
\texttt{"XY"} & $X_1\otimes Y_0$\\
\bottomrule
\end{tabular}
\end{center}

Questa convenzione si applica a \texttt{SparsePauliOp} in Qiskit 2.x ed
è \textbf{invertita} rispetto alla notazione matematica standard (dove
il primo indice è il qubit più a sinistra nel prodotto tensoriale).
Tenerne conto è essenziale nella costruzione di
\texttt{SparsePauliOp.from\_list(...)}.

\section{Sistemi non a qubit: fermioni e bosoni}

Se l'Hamiltoniana descrive elettroni (fermioni) o campi (bosoni), non
nasce direttamente in termini di Pauli, ma in termini di operatori di
creazione e annichilazione $\hat a_i^\dagger,\hat a_j$.

Per i \textbf{sistemi fermionici} (es.\ elettroni in una molecola) si
usano mappature come:
\begin{itemize}[leftmargin=1.3em,itemsep=0.1em]
\item \textbf{Jordan--Wigner}: mappa locale, ma genera stringhe di Pauli
di lunghezza $O(N)$;
\item \textbf{Bravyi--Kitaev}: stringhe mediamente più corte, overhead
minore.
\end{itemize}
Queste trasformazioni preservano le relazioni di anticommutazione e
quindi lo spettro energetico. Per i \textbf{sistemi bosonici} (spazio
di Hilbert infinito-dimensionale) occorre prima troncare lo spazio a un
numero massimo di occupazione, poi mappare i livelli risultanti su
qubit.

\begin{editnote}
Vantaggio dei sistemi spin-1/2 (rilevante per la tesi): ogni spin-1/2 si
mappa \emph{direttamente} su un qubit senza alcun mapping aggiuntivo ---
questo è il punto chiave sottolineato da Crippa et al.\ e dal Prof.\
Chiesa in 6-Applications.
\end{editnote}

\section{Un esempio esplicito: due qubit interagenti}

Consideriamo la seguente matrice hermitiana $4\times4$:
$$H=\begin{pmatrix}1.0&0.0&0.0&0.5\\0.0&-1.0&0.5&0.0\\0.0&0.5&-1.0&0.0\\0.5&0.0&0.0&1.0\end{pmatrix}.$$
Applicando la formula dei coefficienti si ottiene
$$H = 0.5\,(Z\otimes I)+0.5\,(I\otimes Z)+0.5\,(X\otimes X),$$
ovvero, in notazione con pedici fisici, $H=0.5\,Z_1+0.5\,Z_0+0.5\,X_1X_0$.
Il valore atteso su uno stato $|\psi\rangle$ si calcola per linearità:
\begin{keyeq}
\centering
$\langle H\rangle = 0.5\,\langle Z_1\rangle+0.5\,\langle Z_0\rangle+0.5\,\langle X_1X_0\rangle$
\end{keyeq}

\needspace{8\baselineskip}
\section{Il modello di Heisenberg isotropo per il dimero (convenzione 6-Applications)}

\begin{editnote}
Sezione aggiunta per chiarire la convenzione della tesi.
\end{editnote}

In 6-Applications e in Crippa et al., il dimero di Heisenberg è scritto
come
$$H = J_xX_1X_2+J_yY_1Y_2+J_zZ_1Z_2+b(Z_1+Z_2)$$
dove $X_i,Y_i,Z_i$ \textbf{sono già operatori di Pauli} (non operatori
di spin $s_i=\sigma_i/2$). Nel caso isotropo $J_x=J_y=J_z=J$:
\begin{keyeq}
\centering
$H = J(X_1X_2+Y_1Y_2+Z_1Z_2)+b(Z_1+Z_2)$
\end{keyeq}

\begin{editnote}
Attenzione alla convenzione: questo va distinto dall'Hamiltoniana
scritta in termini di operatori di spin $\vec s_i=\vec\sigma_i/2$, che
darebbe $H=2J\,\vec s_1\cdot\vec s_2+\ldots$ con un fattore di scala 4
sulle energie. In 6-Applications la scrittura in Pauli è quella diretta
(nessun fattore $1/4$).
\end{editnote}

L'energia del \textbf{ground state} (singoletto, $b=0$) per questa
convenzione è $E_0=-3J$.

\section{Come si misura una stringa di Pauli su hardware quantistico}

I computer quantistici commerciali (IBM, Rigetti, ecc.) misurano
nativamente solo lungo l'asse $Z$ (base computazionale
$|0\rangle,|1\rangle$). Per misurare un termine diverso si applica una
\textbf{rotazione di base} prima della misura:

\begin{center}
\renewcommand{\arraystretch}{1.25}
\begin{tabular}{@{}cc@{}}
\toprule
\rowcolor{inkblue!10}
\textbf{Operatore da misurare} & \textbf{Porta(e) prima della misura}\\
\midrule
$Z$ & Nessuna (base nativa)\\
$X$ & $H$ (Hadamard)\\
$Y$ & $S^\dagger$ poi $H$\\
\bottomrule
\end{tabular}
\end{center}

La corrispondenza bit $\to$ autovalore è: $0\mapsto+1$, $1\mapsto-1$.

\subsection*{Esempio: misura della stringa $X_0Y_1Z_2$}

\begin{panelbox}{Circuito di misura}
\begin{verbatim}
|psi_0>  --[ H ]------------[ M ]   <- asse X
|psi_1>  --[ S+]--[ H ]-----[ M ]   <- asse Y
|psi_2>  ---------------------[ M ]   <- asse Z (nativo)
\end{verbatim}
\end{panelbox}

Per ogni \emph{shot} (es.\ risultato \texttt{010}):
$$0\to+1,\quad 1\to-1,\quad 0\to+1 \;\Longrightarrow\; (+1)\times(-1)\times(+1)=-1.$$
Il valore atteso $\langle X_0Y_1Z_2\rangle$ è la media su tutti gli
\emph{shot}.

\section{Misura dell'Hamiltoniana del dimero isotropo}

Per $H=J(X_1X_2+Y_1Y_2+Z_1Z_2)$ sono necessari \textbf{tre circuiti di
misura distinti}, poiché le stringhe $XX$, $YY$, $ZZ$ non sono
simultaneamente diagonalizzabili nella stessa base (non commutano tutte
a coppie: ad esempio $[X_1X_2,Y_1Y_2]\neq0$).

\begin{editnote}
Il documento originale usava la formula "non commutano tutte
simultaneamente" senza precisare perché questo implica circuiti
separati. Il motivo è che ciascuna stringa è diagonalizzabile nella
propria base di autostati, ma le tre basi sono diverse, quindi non
esiste un'unica rotazione che diagonalizzi tutte e tre
contemporaneamente.
\end{editnote}

\begin{itemize}[leftmargin=1.3em,itemsep=0.25em]
\item \textbf{Circuito 1} --- misura di $Z_1Z_2$ (base nativa, nessuna
rotazione): $00\to+1,\ 01\to-1,\ 10\to-1,\ 11\to+1$;
\item \textbf{Circuito 2} --- misura di $X_1X_2$ ($H$ su entrambi i
qubit, poi misura): post-processing identico a $ZZ$;
\item \textbf{Circuito 3} --- misura di $Y_1Y_2$ ($S^\dagger H$ su
entrambi i qubit, poi misura): post-processing identico a $ZZ$.
\end{itemize}

\begin{keyeq}
\centering
$\langle H\rangle = J\big(\langle X_1X_2\rangle+\langle Y_1Y_2\rangle+\langle Z_1Z_2\rangle\big)$
\end{keyeq}

\needspace{10\baselineskip}
\section{Verifica analitica: lo stato di singoletto}

Lo stato di singoletto è $|\psi\rangle=\tfrac{1}{\sqrt2}(|01\rangle-|10\rangle)$.

\begin{editnote}
Convenzione ket Qiskit: in $|q_1q_0\rangle$, lo stato $|01\rangle$
significa $q_1=0$, $q_0=1$.
\end{editnote}

\begin{okbox}{Termine $Z_1Z_2$}
$(Z\otimes Z)|\psi\rangle=\frac{1}{\sqrt2}\big[(Z|0\rangle)\otimes(Z|1\rangle)-(Z|1\rangle)\otimes(Z|0\rangle)\big]
=\frac{1}{\sqrt2}\big[-|01\rangle+|10\rangle\big]=-|\psi\rangle
\;\Rightarrow\; \langle Z_1Z_2\rangle=-1$
\end{okbox}
\smallskip

\begin{okbox}{Termine $X_1X_2$}
$(X\otimes X)|\psi\rangle=\frac1{\sqrt2}\big[|10\rangle-|01\rangle\big]=-|\psi\rangle
\;\Rightarrow\; \langle X_1X_2\rangle=-1$
\end{okbox}
\smallskip

\begin{okbox}{Termine $Y_1Y_2$}
$(Y\otimes Y)|\psi\rangle=\frac1{\sqrt2}\big[(i|1\rangle)\otimes(-i|0\rangle)-(-i|0\rangle)\otimes(i|1\rangle)\big]
=\frac1{\sqrt2}\big[|10\rangle-|01\rangle\big]=-|\psi\rangle
\;\Rightarrow\; \langle Y_1Y_2\rangle=-1$
\end{okbox}

\bigskip
\noindent\textbf{Energia totale:}
\begin{keyeq}
\centering
$\langle H\rangle = J(-1-1-1) = -3J$
\end{keyeq}
Per $J>0$ (accoppiamento antiferromagnetico), $-3J$ è il minimo dello
spettro, confermando che il singoletto è lo \textbf{stato fondamentale}.

\begin{editnote}
Raccordo con la convenzione di Crippa: se si parte da
$H=2J\sum_i\vec s_i\cdot\vec s_{i+1}$ con $\vec s_i=\vec\sigma_i/2$, il
termine di scambio diventa $\tfrac{J}2(XX+YY+ZZ)$ e l'energia del
singoletto è $-3J/4$. Il fattore 4 di differenza proviene da
$s_i=\sigma_i/2$. In questo documento si usa la convenzione di
6-Applications dove i coefficienti $J_x,J_y,J_z$ moltiplicano
direttamente gli operatori di Pauli, senza il fattore $1/4$.
\end{editnote}

\needspace{10\baselineskip}
\section{Il vantaggio pratico: perché la decomposizione di Pauli è il pilastro del VQE}

La decomposizione in stringhe di Pauli è essenziale perché:
\begin{enumerate}[leftmargin=1.4em,itemsep=0.25em]
\item \textbf{I computer quantistici non possono misurare una matrice
generica}, ma sanno misurare $X,Y,Z$ su ogni qubit con alta efficienza;
\item \textbf{la linearità} del valor medio permette di calcolare
$\langle H\rangle$ come somma pesata dei valori medi dei singoli
termini;
\item \textbf{la sparsità} delle hamiltoniane fisiche garantisce che il
numero di termini non nulli cresca polinomialmente in $N$ (come
$O(N^k)$ per interazioni a $k$ corpi), rendendo il calcolo scalabile;
\item \textbf{nel caso spin-1/2 $\to$ qubit}, l'Hamiltoniana è
\emph{già} in forma di Pauli senza mapping aggiuntivo --- il vantaggio
chiave sottolineato da Crippa et al.
\end{enumerate}

\end{document}
